
\section{Preliminaries}
\label{sec:preliminary}

%This section defines the problem of autonomous data preparation (ADP), describes the operators supported by \sys, and introduces background on agentic reasoning.

\subsection{Problem Formulation}
\label{subsec:problem_definition}
%This paper studies the problem of \emph{autonomous data preparation} (ADP), which aims to transform and clean tabular data from multiple sources into a target table given only the specification of the target schema.
%We formally define the problem as follows.

\stitle{Source and Target Data.}
We consider a set of \emph{source tables} $\mathcal{S} = \{S_1, \ldots, S_n\}$. Each source table is denoted as $S_i = (\Sigma_i, D_i)$, where $\Sigma_i$ denotes the schema of $S_i$ and $D_i$ denotes the set of tuples conforming to $\Sigma_i$. The schema $\Sigma_i$ captures both table-level and column-level specifications and is defined as $\Sigma_i = (\tau_i, C_i)$. Here, $\tau_i$ represents a textual description that summarizes the semantic meaning of the table, and $C_i = \{c_i^{1}, \ldots, c_i^{m_i}\}$ denotes the set of column specifications. Each column specification $c_i^{j}$ includes column-level metadata such as column name, data type, and semantic description.
In practice, the schema $\Sigma_i$ of a source table may be partially specified or entirely unavailable. In such cases, table- and column-level specifications can be inferred from the raw data in $D_i$. 
For example, each source table $S_i$ in Figure~\ref{fig:introduction_example} consists of a schema $\Sigma_i$ which may be specified by table- and column-level information, and a set $D_i$ of data tuples, \eg (1,``jurassic park'', 101, ``Sci-Fi'').
% For example, consider the ADP case in Figure~\ref{fig:introduction_example}. The schema $\Sigma_i$ of each source table $S_i$ in $\mathcal{S}$ is partially specified with table name (\eg ``movies'') and column name (\eg ``id''). The raw 

% \fanj{For example, \ldots}

Similarly, the \emph{target table} is denoted as $T=(\Sigma^{\ast}, D^{\ast})$, where $\Sigma^\ast = (\tau^\ast, C^\ast)$ specifies the target schema and $D^\ast$ denotes the target tuples to be produced. 
Note that, in autonomous data preparation, only the target schema $\Sigma^\ast$ is given as input, while the target tuples $D^\ast$ are unknown and must be generated by transforming and cleaning the source data.
For example, in Figure~\ref{fig:introduction_example}, the target schema $\Sigma^{\ast}$ is specified by a table-level description (\eg ``\emph{I need a table containing \ldots}'') and a set of column specifications, each including a column name (\eg ``\emph{director\_name}'') and a corresponding semantic description (\eg ``\emph{The name of the director.}'').
% \fanj{For example, \ldots}

%Let $T$ be a source table consisting of a schema (columns) and instances (rows). Let $\mathbb{O}$ be the set of operators. We formulate the problem of Autonomous Data Preparation (ADP) as below.
%
%\begin{definition}
%    Given (1) a collection of \textit{source tables} $\mathcal{D}=\{T_1,\dots,T_M\}$ where $T$ denotes a source table consisting of a table name, a schema (columns) and instances (rows), and (2) \textit{a target specification} $\mathcal{S}=(\mathcal{S}_{table}, \mathcal{S}_{column})$ where $\mathcal{S}_{table}$ denotes a table-level summary of the target table and $\mathcal{S}_{column}$ denotes a structured constraint specifying the desired output schema. 
%    Synthesize a pipeline of parameterized operators $\Pi=[o_1, \dots, o_i]$, with $type(o_i) \in \mathbb{O}$, such that applying each operator $o_i$ successively on $\mathcal{D}$ with the operator execution engin $\mathcal{F}_{op}$ to produce a generated table $\hat{T}$ that aligns with the target specification $\mathcal{S}$.
%\end{definition}

\stitle{Data Preparation Pipeline.}
We model the process of data preparation as a composition of \emph{operators} applied to source tables.
Specifically, an operator represents a primitive data transformation. 
%We consider a finite set of \emph{operator types}, where each type specifies the transformation semantics and the corresponding parameters. 
%
Formally, an operator $o$ is characterized by its type $\kappa$ and parameters $\theta$, and defines a function $o_{\theta}: \mathcal{T} \rightarrow \mathcal{T}$, where $\mathcal{T}$ denotes a set of tables, and each table is of the form $(\Sigma, D)$ as defined previously. Given an input table set $\mathcal{T}_{\tt in}$, the operator produces an output table set, \ie $\mathcal{T}_{\tt out}=o_{\theta}(\mathcal{T}_{\tt in})$.
Note that different operator types capture different classes of data transformations, such as filtering, joining, aggregation, and cleaning. The current operator types supported by \sys will be introduced in Section~\ref{subsec:operator_space}.

We then define a data preparation pipeline (or pipeline for short) as an ordered sequence of operators applied iteratively to transform a set of source tables $\mathcal{S}$ into a target table $\hat{T}$. We use the symbol ``$\to$'' to denote the left-to-right application order between adjacent operators, \ie the output produced by an operator is fed as the input to the operator to its right. Formally, a pipeline $\mathcal{P}$ is defined as an ordered composition of $k$ operators over $\mathcal{S}$, \ie
%
\begin{equation} 
\label{eq:pipeline_definition}
    \mathcal{P} = (\mathcal{S}, o^{(1)}_{\theta_1} \to o^{(2)}_{\theta_2} \to \cdots \to o^{(k)}_{\theta_k}).
\end{equation}
%
Given a set of source tables $\mathcal{S}$, the execution of pipeline $\mathcal{P}$ produces a sequence of intermediate table sets, \ie 
%
$
    \mathcal{T}_{0} = \mathcal{S}, ~\mathcal{T}_i =  o^{(i)}_{\theta_i}(\mathcal{T}_{i-1}), i=1,\ldots,k, 
$
%
where $\mathcal{T}_k$ denotes the final output of the pipeline, which contains only the target table $\hat{T}$, \ie $\mathcal{T}_k=\{\hat{T}\}$.
%
%\add{
Consider operator \texttt{Deduplicate (movies, [id], first)} as an example. 
Its operator type $\kappa$ is \texttt{Deduplicate} and its parameters $\theta$ refer to (\texttt{\textit{movies}, [id], first}). Executing this operator on source tables produces intermediate tables consisting of a cleaned \textit{movie}, while tables \textit{ratings} and \textit{directors} remain unchanged.
% For example, $\mathcal{T}_0$ is initialized with source tables illustrated in the left part of Figure~\ref{fig:introduction_example}. By executing operator, \eg \texttt{Deduplicate} that removes duplicated record in table ``movies'' based on id values, we transform $\mathcal{T}_0$ to $\mathcal{T}_1$ which consists of updated table ``movies''
%}
% \fanj{For Example, \ldots}

\stitle{Problem Statement.}
Given a set of source tables $\mathcal{S} = \{S_1, \ldots, S_n\}$, a target schema $\Sigma^{\ast}$, and a set of available operator types $\mathcal{O}$, the autonomous data preparation (ADP) problem is to find an optimal data preparation pipeline $\mathcal{P}$ that transforms $\mathcal{S}$ into a target table $\hat{T} = (\Sigma^{\ast}, \hat{D})$, where $\hat{D}$ denotes the tuples produced by $\mathcal{P}$ that satisfy the target schema $\Sigma^{\ast}$.
We assume a ground-truth target table $T^\ast=(\Sigma^\ast, D^\ast)$, which is used \emph{only} for evaluation. Thus, the optimality of a pipeline is defined by the quality of $\hat{D}$ with respect to $D^\ast$ under task-specific evaluation metrics, which will be introduced in the experimental setup (Section~\ref{subsec:exp-setup}).
%
\begin{example}

Figure~\ref{fig:introduction_example} provides an example of an ADP task with three source tables: \textit{movies}, \textit{ratings}, and \textit{directors}. The goal is to generate a target table that conforms to a given target schema. The target schema describes both the overall semantics of the table and its desired structure. Specifically, it provides a table-level summary of the intended content and specifies column-level requirements, including column names (\eg \textit{genre}), semantic meanings (\eg movie genre category), and data constraints (\eg single-valued values).
To achieve this goal, we construct a data preparation pipeline composed of multiple operators, such as \texttt{Deduplicate} and \texttt{Join}. Each operator is applied to the current set of tables and produces updated intermediate tables. For example, applying \texttt{Join} merges the \textit{movies} and \textit{directors} tables on movie titles, yielding a table named \textit{movies\_directors\_join} that integrates movie data with director names, while keeping the \textit{ratings} table unchanged.
After executing all operators in the pipeline, the process produces the target table, shown on the right of Figure~\ref{fig:introduction_example}. 
%    Consider the ADP task $(\mathcal{D}, \mathcal{S})$ in Figure~\ref{fig:introduction_example}.
%    As shown in the left part, $\mathcal{D}$ contains three source tables which are named as \textit{movies}, \textit{ratings} and \textit{directors} respectively.
%    $\mathcal{S}=(\mathcal{S}_{table}, \mathcal{S}_{column})$ describes what does the target table look like. Specifically, $\mathcal{S}_{table}$ provides table-level summary, which is \textit{``The table contains''} and $\mathcal{S}_{column}$ provides description for the schema in the table which specifies the column name (\eg ``genre''), the description for the semantic meanings (\eg ``movie genre category'') and the requirement for the data (\eg ``A single'').
%    It requires to generate a ground-truth operator pipeline $\Pi$ consisting of operators including \texttt{Deduplicate}, \texttt{Join}, etc.
%    These operators will be executed to modify source tables $\mathcal{D}$ by updating one or more tables in $\mathcal{D}$.
%    For example, we conduct \texttt{Deduplicate} to filter out duplicated rows in table \textit{movies} and update the source tables $\mathcal{D}$ with original table \textit{ratings}, original table \textit{directors} and a cleaned table \textit{movies}.
%    Finally, after executing all operators in the ground-truth operator pipeline, it outputs a table $\hat{T}$ shown on the right of Figure~\ref{fig:introduction_example}. The generated table contains data from source tables which aligns with the target specification $\mathcal{S}$.
\end{example}

%\stitle{Remarks.}
%The design of target schema $\Sigma^\ast$ removes the need for access to ground-truth target data, which is often unavailable in real-world settings. This distinguishes our problem from prior work~\cite{he2018transform, gulwani2012spreadsheet, jin2017foofah, singh2016blinkfill} that assumes access to target data examples.

%The design of the target specification $\mathcal{S}$ ensures that the method does not need the access to the data instances of ground-truth target table, which is difficulty to get in real application scenarios. This is the essential difference between our problem and the one that the previous work~\cite{he2018transform, gulwani2012spreadsheet, jin2017foofah, singh2016blinkfill} focused on.

\subsection{Supported Data Preparation Operators}
\label{subsec:operator_space}
\sys supports a broad set of data preparation operators and is designed to be extensible, allowing new operator types to be incorporated as needed. In total, we consider 31 operators, organized into eight categories that cover the major transformation patterns in data preparation pipelines.
% , including data cleaning, normalization, structural reshaping, table combination, and aggregation. 
Due to space constraints, we provide descriptions of \textbf{typical operators} in each category, with complete specifications available in~\cite{technical_report}.

% \begin{table}[t]
  \centering
  \caption{\bf Overview of the Operator Space in \sys.}
  \vspace{-1em}
  \label{tbl:operator_space}
  \resizebox{1.0\columnwidth}{!}{
    \begin{tabular}{|c||l|}
      \hline
      \textbf{Category} & \textbf{Operators} \\
      \hline \hline

      % --- Cell-level ---
      \textbf{Cell-level}
      & \texttt{ValueTransform}, \texttt{StandardizeDatetime}, \texttt{CastType} \\
      \hline \hline

      % --- Row-level ---
      \multirow{3}{*}{\textbf{Row-level}}
      & \texttt{Filter}, \texttt{Sort}, \texttt{TopK}, \texttt{Deduplicate} \\
      & \texttt{DropNA}, \texttt{MissingValueImputation}, \texttt{ErrorDetection} \\
      & \texttt{OutlierDetection} \\
      \hline \hline

      % --- Column-level ---
      \multirow{2}{*}{\textbf{Column-level}}
      & \texttt{SplitColumn}, \texttt{RenameColumn}, \texttt{AddNewColumn}, \texttt{DropColumn} \\
      & \texttt{Concatenate}, \texttt{SelectColumn}, \texttt{Subtitle} \\
      \hline \hline

      % --- Table-level ---
      \multirow{2}{*}{\textbf{Table-level}}
      & \texttt{Pivot}, \texttt{Stack}, \texttt{WideToLong}, \texttt{Transpose}, \texttt{Explode} \\
      & \texttt{Join}, \texttt{Union}, \texttt{Append} \\
      \hline \hline

      % --- Aggregation ---
      \textbf{Aggregation}
      & \texttt{GroupBy}, \texttt{Count}, \texttt{CalculateStatistic} \\
      \hline \hline

      % --- Escape Hatch ---
      \textbf{Escape Hatch}
      & \texttt{ExeCode} \\
      \hline

      % If you must include Terminate, uncomment the following:
      % \hline \hline
      % \textbf{Control} & \texttt{Terminate} \\
      % \hline

    \end{tabular}
  }
  \vspace{-0.5em}
\end{table}




%\add{
%\begin{table*}[t]
  \centering
  \caption{\bf Taxonomy of supported data preparation operators by functional type.}
  \vspace{-0.6em}
  \label{tbl:operator_space_type}

  \small
  \setlength{\tabcolsep}{6pt}
  \renewcommand{\arraystretch}{1.12}

  \resizebox{\textwidth}{!}{
    \begin{tabular}{|l||l|}
      \hline
      \textbf{Operator Type} & \textbf{Operators} \\
      \hline \hline

      \textbf{Data Cleaning} &
      \texttt{MissingValueImputation},
      \texttt{Deduplicate},
      \texttt{ErrorDetection},
      \texttt{OutlierDetection},
      \texttt{DropNA} \\
      \hline

      \textbf{Value Normalization} &
      \texttt{ValueTransform},
      \texttt{StandardizeDatetime},
      \texttt{CastType} \\
      \hline

      \textbf{Schema Editing} &
      \texttt{RenameColumn},
      \texttt{SplitColumn},
      \texttt{AddNewColumn},
      \texttt{DropColumn},
      \texttt{Concatenate},
      \texttt{SelectColumn},
      \texttt{Subtitle} \\
      \hline

      \textbf{Row Selection} &
      \texttt{Filter},
      \texttt{Sort},
      \texttt{TopK} \\
      \hline

      \textbf{Aggregation} &
      \texttt{GroupBy},
      \texttt{Count},
      \texttt{CalculateStatistic} \\
      \hline

      \textbf{Table Combination} &
      \texttt{Join},
      \texttt{Union},
      \texttt{Append} \\
      \hline

      \textbf{Table Reshaping} &
      \texttt{Pivot},
      \texttt{Stack},
      \texttt{WideToLong},
      \texttt{Transpose},
      \texttt{Explode} \\
      \hline

      \textbf{Program Synthesis} &
      \texttt{ExeCode} \\
      \hline
    \end{tabular}
  }
\end{table*}

%}

\subsubsection{Data Cleaning}
Data cleaning operators improve data quality by detecting, fixing, or removing erroneous, missing, or duplicate records, and are commonly applied in data preparation pipelines.
%
\begin{itemize}[leftmargin=*]
    % \item \texttt{DropNA(table, subset, how)} removes rows with missing values in specified columns, with \texttt{subset} specifying the columns and \texttt{how} (``any'' or ``all'') controlling the removal criterion.
    \item \texttt{MissingValueImputation(table, column, mode)} imputes missing values in a specified column using a statistical strategy, with \texttt{mode} selecting the method (\eg mean, median, or mode).
    \item \texttt{Deduplicate(table, subset, keep)} removes duplicate rows based on specified columns, retaining either the first or last occurrence as determined by \texttt{keep}.
    % \item \texttt{ErrorDetection(table, column, func)} identifies invalid records in a specified column based on a user-defined function.
    % \item \texttt{OutlierDetection(table, column, action)} detects statistical outliers in a specified numerical column, with \texttt{action} determining whether outliers are removed or flagged.
\end{itemize}

\subsubsection{Value Normalization}
Value normalization operators standardize field values and data types to ensure semantic consistency.
%
\begin{itemize}[leftmargin=*]
    \item \texttt{ValueTransform(table, column, func)} applies a user-defined transformation function to values in a specified column to standardize their format.
    \item \texttt{StandardizeDatetime(table, column, format)} converts values in a specified column to a user-defined datetime format.
    % \item \texttt{CastType(table, column, dtype)} casts values in a specified column to a target data type defined by \texttt{dtype}.
\end{itemize}

\subsubsection{Schema Editing}
Schema editing operators modify table structure and column semantics, enabling column selection, renaming, derivation, and restructuring to align raw data with the target schema.
%
\begin{itemize}[leftmargin=*]
    \item \texttt{RenameColumn(table, rename\_map)} renames columns according to a specified mapping from original to target names.
    \item \texttt{AddNewColumn(table, name, func)} adds a new column computed from existing data using a row-wise function.
    % \item \texttt{DropColumn(table, columns)} removes specified columns.
    % \item \texttt{SplitColumn(table, source, target, func)} splits a column into multiple columns using a user-defined splitting function.
    % \item \texttt{Concatenate(table, columns, target, func)} combines multiple columns into a single column using a formatting function.
    % \item \texttt{SelectColumn(table, columns)} retains only the specified columns.
    % \item \texttt{Subtitle(table, title, target\_col)} adds a new column populated with a constant title or subtitle value.
\end{itemize}

\subsubsection{Row Selection}
Row selection operators manipulate tables by filtering, ordering, or selecting subsets of rows.
%
\begin{itemize}[leftmargin=*]
    \item \texttt{Filter(table, func)} selects rows satisfying a given predicate.
    \item \texttt{Sort(table, by, ascending)} orders rows in a table based on specified columns, with \texttt{ascending} controlling the direction.
    \item \texttt{TopK(table, k)} retains the top-$k$ rows of the table.
\end{itemize}

\subsubsection{Aggregation}
Aggregation operators summarize data by grouping records and computing statistics, transforming row-level data into compact, analysis-ready representations.
%
\begin{itemize}[leftmargin=*]
    \item \texttt{GroupBy(table, by, agg)} groups rows by specified columns and applies aggregation functions defined by \texttt{agg}.
    \item \texttt{Count(table)} returns the total number of rows in the table. We treat \texttt{Count} as a primitive aggregation operator for convenience, though it can be viewed as a special case of \texttt{GroupBy} without grouping keys.
    % \item \texttt{CalculateStatistic(table, stat, func)} computes a global statistic over the table using a user-defined function, returning a scalar value that summarizes a table-level property.
\end{itemize}

\subsubsection{Table Combination}
Table combination operators integrate data from multiple tables into a unified representation, enabling multi-source data preparation pipelines.
%
\begin{itemize}[leftmargin=*]
    \item \texttt{Join(left, right, on, how)} combines two tables based on specified key columns using a join type defined by \texttt{how}.
    \item \texttt{Union(tables, how)} vertically combines multiple tables with compatible schemas, with \texttt{how} determining whether duplicate rows are retained or removed.
    % \item \texttt{Append(table, other)} appends one table to another by vertically concatenating their rows, without performing duplicate elimination as in \texttt{Union}.
\end{itemize}

\subsubsection{Table Reshaping}
Table reshaping operators reorganize the structural layout of tables, such as converting between wide and long formats or expanding nested fields.
%
\begin{itemize}[leftmargin=*]
    \item \texttt{Pivot(table, index, columns, values, aggfunc)} reshapes a table from long to wide format by aggregating values at the intersection of specified index and column identifiers.
    \item \texttt{Stack(table, id\_vars, value\_vars)} reshapes a table to long format by collapsing multiple columns into key–value pairs.
    % \item \texttt{WideToLong(table, stubnames, i, j)} reshapes wide-format data into long format by parsing column naming patterns.
    % \item \texttt{Transpose(table)} swaps the rows and columns of the table.
    % \item \texttt{Explode(table, column)} expands a column containing list-valued entries into multiple rows.
\end{itemize}

\subsubsection{Program Synthesis}
The program synthesis operator provides a general mechanism for handling transformations that cannot be expressed using the above predefined operators.
\begin{itemize}[leftmargin=*]
    \item \texttt{ExeCode(tables, target, func)} generates a target table from input tables by executing Python code synthesized by an LLM, enabling transformations that cannot be expressed using the above predefined operators.
\end{itemize}
